\documentclass[aspectratio=169]{beamer}

\usetheme{Madrid}
\usecolortheme{default}
\setbeamertemplate{navigation symbols}{}

\usepackage{amsmath}
\usepackage{amssymb}
\usepackage{booktabs}
\usepackage{tikz}
\usetikzlibrary{arrows.meta, positioning, calc}
\usepackage{tikz-3dplot}
\usepackage{algorithm}
\usepackage[noend]{algpseudocode}
\usepackage{xspace}
\usepackage[all]{xy}
\usepackage{svg}
\usepackage{siunitx}
\sisetup{detect-weight=true,group-four-digits=true}
\usepackage{subcaption}
\usepackage[absolute,overlay]{textpos}
\usepackage{colortbl}
\usepackage[nopostdot,nogroupskip,nonumberlist]{glossaries}

% prelude.tex — shared macros for paper and presentation
% Usage: % prelude.tex — shared macros for paper and presentation
% Usage: % prelude.tex — shared macros for paper and presentation
% Usage: \input{prelude} after \documentclass and package loading

% ── Glossary / acronyms ─────────────────────────────────────────────────────
\newacronym{sat}{SAT}{\emph{Boolean Satisfiability}}
\newacronym{maxSat}{maxSAT}{\emph{Maximum Satisfiability}}
\newacronym{pb}{PB}{\emph{Pseudo-Boolean}}
\newacronym{cp}{CP}{\emph{Constraint Programming}}
\newacronym{tsp}{TSP}{\emph{Travelling Salesperson Problem}}
\newacronym{mdd}{MDD}{\emph{Multivalued Decision Diagram}}
\newacronym{bdd}{BDD}{\emph{Binary Decision Diagram}}
\newacronym{csp}{CSP}{\emph{Constraint Satisfaction Problem}}
\newacronym{cop}{COP}{\emph{Constrained Optimisation Problem}}
\newacronym{mip}{MIP}{\emph{Mixed Integer Programming}}
\newacronym{ilp}{ILP}{\emph{Integer Linear Programming}}
\newacronym{lp}{LP}{\emph{Linear Programming}}
\newacronym{lbbd}{LBBD}{\emph{Logic-Based Benders Decomposition}}
\newacronym{mus}{MUS}{\emph{Minimal Unsatisfiable Subset}}

\newcommand{\milp}{\gls{ilp}\xspace}
\newcommand{\relaxation}{\acrshort{lp} relaxation\xspace}

% ── Text shorthands ──────────────────────────────────────────────────────────
\newcommand{\dquote}[1]{``#1''}
\newcommand{\quoted}[1]{`#1'}
\newcommand{\ie}{i.e.\xspace}
\newcommand{\eg}{e.g.\xspace}
\newcommand{\Eg}{E.g.\xspace}

% ── Author annotation commands (silenced by default) ─────────────────────────
\newcommand{\question}[1]{}
\newcommand{\pjs}[1]{}
\newcommand{\tias}[1]{}
\newcommand{\hb}[1]{}
\newcommand{\wout}[1]{}
\newcommand{\scrap}[1]{}
\newcommand{\ignore}[1]{}

\newcommand{\red}[1]{\textcolor{blue}{#1}}

% ── Math notation ────────────────────────────────────────────────────────────
\newcommand{\brackets}[1]{\ensuremath{[\![#1]\!]}}
\newcommand{\iverson}[1]{\ensuremath{\langle #1 \rangle}}
\newcommand{\beq}[2]{\brackets{#1 = #2}}
\newcommand{\bevarlookup}[1]{\ifnum9<1#1 \ifcase#1\or x\or y\or z\fi\else x_{#1}\fi}
\newcommand{\x}[1]{\bevarlookup{#1}}
\newcommand{\be}[2]{\beq{\bevarlookup{#1}}{#2}}
\newcommand{\besml}[2]{\scriptstyle\be{#1}{#2}}

\newcommand{\etAl}[1]{#1~\emph{et al.}}
\newcommand{\citeNoun}[2]{\etAl{#1}~\cite{#2}}

\DeclareMathOperator*{\argmax}{arg\,max}
\DeclareMathOperator*{\argmin}{arg\,min}

\newcommand{\assignment}{\ensuremath{\hat{v}}}
\newcommand{\true}{\mathit{true}}
\newcommand{\false}{\mathit{false}}
\newcommand{\Integers}{\mathbb{Z}}
\newcommand{\Reals}{\mathbb{R}}
\newcommand{\Bools}{\mathbb{B}}
\renewcommand{\emptyset}{\varnothing}

\newcommand{\vars}[1]{\ensuremath{vars(#1)}}
\newcommand{\scope}[1]{\ensuremath{scope(#1)}}
\newcommand{\cons}[1]{\ensuremath{cons(#1)}}
\newcommand{\bounds}[1]{\ensuremath{bounds(#1)}}

\newcommand{\interval}[2]{#1..#2}
\newcommand{\tuple}[1]{\langle #1 \rangle}
\newcommand{\sequence}[1]{[#1]}
\newcommand{\set}[1]{\{#1\}}
\newcommand{\setComp}[2]{\set{#1 ~|~ #2}}
\newcommand{\sequenceComp}[2]{\sequence{#1 ~|~ #2}}
\newcommand{\dom}[1]{D(#1)}
\newcommand{\pp}{+\!\!+}

\newcommand{\rhs}{s}

% ── Table-specific notation ──────────────────────────────────────────────────
\newcommand{\cols}{n}
\newcommand{\rows}{m}
\newcommand{\row}{r}
\newcommand{\col}{j}

\newcommand{\llinex}[1]{Line \ref{#1}}
\newcommand{\linex}[1]{line \ref{#1}}
\newcommand{\lines}[2]{lines \ref{#1}--\ref{#2}}

\newcommand{\ttable}{\texttt{table}\xspace}
\newcommand{\btable}{T^{\Bools}}
\newcommand{\btableTransposed}{\hat{T}^{\Bools}}
\newcommand{\transpose}[1]{\hat{#1}}

\newcommand{\nolbbd}[2]{#2}

% ── Configuration / approach labels ──────────────────────────────────────────
\newcommand{\textconfig}[1]{\textsc{#1}}
\newcommand{\baseGleb}{\textconfig{Int}\xspace}
\newcommand{\baseBool}{\textconfig{Bool}\xspace}

\newcommand{\baseMdd}{\textconfig{MDD}\xspace}
\newcommand{\baseMddInput}{\textconfig{\baseMdd-input}\xspace}
\newcommand{\baseMddDomIncr}{\textconfig{\baseMdd-dom}\xspace}
\newcommand{\baseMddGreedy}{\textconfig{\baseMdd-greedy}\xspace}

\newcommand{\lazyGenerate}{\textconfig{Generate}\xspace}
\newcommand{\lazyNone}{\textconfig{Generate}\xspace}
\newcommand{\lazyFractional}{\textconfig{Fractional}\xspace}
\newcommand{\lazyNegatives}{\textconfig{Negatives}\xspace}
\newcommand{\lazyCutlift}{\textconfig{Cutlift}\xspace}
\newcommand{\lazyShrink}{\textconfig{Shrink}\xspace}
\newcommand{\lazyCutoff}{\textconfig{Hybrid}\xspace}
\newcommand{\lazyCutoffT}[1]{\textconfig{Hybrid (#1)}\xspace}
\newcommand{\lazyAll}{\textconfig{All}\xspace}

% ── xy-pic node macros ───────────────────────────────────────────────────────
\newcommand{\xyo}[1]{*++[o][F-]{#1}}
\newcommand{\xyob}[1]{*++[o][F=]{#1}}
\newcommand{\xyoy}[1]{*++[o][F=]{#1}}
\newcommand{\xyoo}[1]{*++[o][F=]{#1}}

% ── MDD node labels (shared by \mddDiagram*, \mddFlowA--D) ────────────────────
\newcommand{\mddNodeRoot}{[]}
\newcommand{\mddNodeXtwo}{[2]}
\newcommand{\mddNodeXone}{[1]}
\newcommand{\mddNodeXfour}{[4]}
\newcommand{\mddNodeXthree}{[3]}
\newcommand{\mddNodeXtwoYone}{[2,1]}
\newcommand{\mddNodeMerge}{\scriptscriptstyle[1,2] \scriptscriptstyle[4,3]}
\newcommand{\mddNodeXthreeYtwo}{[3,2]}
\newcommand{\mddNodeSnk}{snk}

% ── Algorithm macros ─────────────────────────────────────────────────────────
\newcommand{\codeBlockDelim}{}
\newcommand{\tableCon}{\ensuremath{\texttt{table}(\hat{b}, \btable)}}
\newcommand{\generateInitArgs}{\ensuremath{\assignment, \tableCon}}
\newcommand{\generateArgs}{\ensuremath{\generateInitArgs, X, \rhs, R, V}}
\newcommand{\validCut}{\sum_{\col \in X} a_\col b_\col \leq \rhs}

% ── Cut lifting macros ──────────────────────────────────────────────────────
\newcommand{\Cols}{\interval{1}{\cols}}
\newcommand{\Rows}{\interval{1}{\rows}}
\newcommand{\colLift}{\col'}
\newcommand{\Rtight}{\ensuremath{R_{\text{tight}}}}
\newcommand{\Ntight}{\ensuremath{N_{\text{tight}}}}
\newcommand{\Xtight}{\ensuremath{X'}}
\newcommand{\cutliftThm}{%
Lifting is correct.
Suppose $\btable \models \sum_{\col \in X} a_\col b_\col \leq \rhs$  and $\btable \models b_{\colLift} \rightarrow \sum_{\col \in X} a_\col b_\col \leq \rhs-a_{\colLift}$ then
$\btable \models \sum_{\col \in X \cup \set{\colLift}} a_\col b_\col \leq \rhs$.%
}

% ── Negatives macros ─────────────────────────────────────────────────────────
\newcommand{\negate}[1]{\overline{#1}}
\newcommand{\partSize}[1]{\hat{P}_{#1}}
\newcommand{\unencoded}{B}

% ── Experimental evaluation ──────────────────────────────────────────────────
\newcommand{\timeout}{600\xspace}
\newcommand{\memout}{8\xspace}
\newcommand{\unk}{\texttt{UNK}\xspace}
\newcommand{\mem}{\texttt{MEM}\xspace}
\newcommand{\pst}{\texttt{POST}\xspace}
\newcommand{\sat}{\texttt{FEAS}\xspace}
\newcommand{\uns}{\texttt{UNS}\xspace}
\newcommand{\sol}{\texttt{SOLVED}\xspace}
\newcommand{\opt}{\texttt{OPT}\xspace}
\newcommand{\XCSPYears}{XCSP3'22\textendash25\xspace}
\newcommand{\thousands}{$10^3$}

% ── Figures ───────────────────────────────────────────────────────────────────
% Locally thicken xy-pic edges/frames within a group (restored at the closing brace)
\makeatletter
\newcommand{\xyThick}{\begingroup\xylinethick@=\dimexpr2\xylinethick@\relax}
\newcommand{\xyThickEnd}{\endgroup}
\makeatother
% MDD diagram for the running example (parameterised by xy spacing)
\newcommand{\mddDiagram}[2]{% #1 = @C (column sep), #2 = @R (row sep)
\xyThick
\xymatrix@C=#1@R=#2{
   \xyo{\mddNodeRoot}
       \ar[r]^{\be{1}{2}}
       \ar[rd]^{\be{1}{1}}
       \ar[rdd]^>>>>>>{\be{1}{4}}
       \ar[rddd]_{\be{1}{3}}
       & \xyo{\mddNodeXtwo} \ar[r]^{\be{2}{1}}
       & \xyo{\mddNodeXtwoYone} \ar[r]^{\be{3}{1}} \ar@/_2pc/[r]^{\be{3}{2}}
       & \xyoo{\mddNodeSnk} \\
   & \xyo{\mddNodeXone} \ar[rd]^{\be{2}{2}} &  \\
   & \xyo{\mddNodeXfour} \ar[r]_{\be{2}{3}}
       & \xyo{\mddNodeMerge} \ar[ruu]^{\be{3}{3}} \\
   & \xyo{\mddNodeXthree} \ar[r]^{\be{2}{2}}
       & \xyo{\mddNodeXthreeYtwo} \ar[ruuu]_{\be{3}{2}} \\
}%
\xyThickEnd
}
\newcommand{\mddDiagramHighlight}[2]{% same but with path 1,2,3 highlighted
\xyThick
\xymatrix@C=#1@R=#2{
   \xyo{\mddNodeRoot}
       \ar[r]^{\be{1}{2}}
       \ar@[blue][rd]^{\textcolor{blue}{\be{1}{1}}}
       \ar[rdd]^>>>>>>{\be{1}{4}}
       \ar[rddd]_{\be{1}{3}}
       & \xyo{\mddNodeXtwo} \ar[r]^{\be{2}{1}}
       & \xyo{\mddNodeXtwoYone} \ar[r]^{\be{3}{1}} \ar@/_2pc/[r]^{\be{3}{2}}
       & \xyoo{\mddNodeSnk} \\
   & \xyo{\textcolor{blue}{\mddNodeXone}} \ar@[blue][rd]^{\textcolor{blue}{\be{2}{2}}} &  \\
   & \xyo{\mddNodeXfour} \ar[r]_{\be{2}{3}}
       & \xyo{\textcolor{blue}{\mddNodeMerge}} \ar@[blue][ruu]^{\textcolor{blue}{\be{3}{3}}} \\
   & \xyo{\mddNodeXthree} \ar[r]^{\be{2}{2}}
       & \xyo{\mddNodeXthreeYtwo} \ar[ruuu]_{\be{3}{2}} \\
}%
\xyThickEnd
}
\newcommand{\mddDiagramHighlightBoth}[2]{% paths 1,2,3 (blue) and 4,3,3 (red)
\xyThick
\xymatrix@C=#1@R=#2{
   \xyo{\mddNodeRoot}
       \ar[r]^{\be{1}{2}}
       \ar@[blue][rd]^{\textcolor{blue}{\be{1}{1}}}
       \ar@[red][rdd]^>>>>>>{\textcolor{red}{\be{1}{4}}}
       \ar[rddd]_{\be{1}{3}}
       & \xyo{\mddNodeXtwo} \ar[r]^{\be{2}{1}}
       & \xyo{\mddNodeXtwoYone} \ar[r]^{\be{3}{1}} \ar@/_2pc/[r]^{\be{3}{2}}
       & \xyoo{\mddNodeSnk} \\
   & \xyo{\textcolor{blue}{\mddNodeXone}} \ar@[blue][rd]^{\textcolor{blue}{\be{2}{2}}} &  \\
   & \xyo{\textcolor{red}{\mddNodeXfour}} \ar@[red][r]_{\textcolor{red}{\be{2}{3}}}
       & \xyo{\textcolor{violet}{\mddNodeMerge}} \ar@[violet][ruu]^{\textcolor{violet}{\be{3}{3}}} \\
   & \xyo{\mddNodeXthree} \ar[r]^{\be{2}{2}}
       & \xyo{\mddNodeXthreeYtwo} \ar[ruuu]_{\be{3}{2}} \\
}%
\xyThickEnd
}

% ── Misc ─────────────────────────────────────────────────────────────────────
\newcommand{\omitVarJ}{[x_1, \ldots, x_{j-1}, x_{j+1}, \ldots, x_{n-1}, x_n]}
 after \documentclass and package loading

% ── Glossary / acronyms ─────────────────────────────────────────────────────
\newacronym{sat}{SAT}{\emph{Boolean Satisfiability}}
\newacronym{maxSat}{maxSAT}{\emph{Maximum Satisfiability}}
\newacronym{pb}{PB}{\emph{Pseudo-Boolean}}
\newacronym{cp}{CP}{\emph{Constraint Programming}}
\newacronym{tsp}{TSP}{\emph{Travelling Salesperson Problem}}
\newacronym{mdd}{MDD}{\emph{Multivalued Decision Diagram}}
\newacronym{bdd}{BDD}{\emph{Binary Decision Diagram}}
\newacronym{csp}{CSP}{\emph{Constraint Satisfaction Problem}}
\newacronym{cop}{COP}{\emph{Constrained Optimisation Problem}}
\newacronym{mip}{MIP}{\emph{Mixed Integer Programming}}
\newacronym{ilp}{ILP}{\emph{Integer Linear Programming}}
\newacronym{lp}{LP}{\emph{Linear Programming}}
\newacronym{lbbd}{LBBD}{\emph{Logic-Based Benders Decomposition}}
\newacronym{mus}{MUS}{\emph{Minimal Unsatisfiable Subset}}

\newcommand{\milp}{\gls{ilp}\xspace}
\newcommand{\relaxation}{\acrshort{lp} relaxation\xspace}

% ── Text shorthands ──────────────────────────────────────────────────────────
\newcommand{\dquote}[1]{``#1''}
\newcommand{\quoted}[1]{`#1'}
\newcommand{\ie}{i.e.\xspace}
\newcommand{\eg}{e.g.\xspace}
\newcommand{\Eg}{E.g.\xspace}

% ── Author annotation commands (silenced by default) ─────────────────────────
\newcommand{\question}[1]{}
\newcommand{\pjs}[1]{}
\newcommand{\tias}[1]{}
\newcommand{\hb}[1]{}
\newcommand{\wout}[1]{}
\newcommand{\scrap}[1]{}
\newcommand{\ignore}[1]{}

\newcommand{\red}[1]{\textcolor{blue}{#1}}

% ── Math notation ────────────────────────────────────────────────────────────
\newcommand{\brackets}[1]{\ensuremath{[\![#1]\!]}}
\newcommand{\iverson}[1]{\ensuremath{\langle #1 \rangle}}
\newcommand{\beq}[2]{\brackets{#1 = #2}}
\newcommand{\bevarlookup}[1]{\ifnum9<1#1 \ifcase#1\or x\or y\or z\fi\else x_{#1}\fi}
\newcommand{\x}[1]{\bevarlookup{#1}}
\newcommand{\be}[2]{\beq{\bevarlookup{#1}}{#2}}
\newcommand{\besml}[2]{\scriptstyle\be{#1}{#2}}

\newcommand{\etAl}[1]{#1~\emph{et al.}}
\newcommand{\citeNoun}[2]{\etAl{#1}~\cite{#2}}

\DeclareMathOperator*{\argmax}{arg\,max}
\DeclareMathOperator*{\argmin}{arg\,min}

\newcommand{\assignment}{\ensuremath{\hat{v}}}
\newcommand{\true}{\mathit{true}}
\newcommand{\false}{\mathit{false}}
\newcommand{\Integers}{\mathbb{Z}}
\newcommand{\Reals}{\mathbb{R}}
\newcommand{\Bools}{\mathbb{B}}
\renewcommand{\emptyset}{\varnothing}

\newcommand{\vars}[1]{\ensuremath{vars(#1)}}
\newcommand{\scope}[1]{\ensuremath{scope(#1)}}
\newcommand{\cons}[1]{\ensuremath{cons(#1)}}
\newcommand{\bounds}[1]{\ensuremath{bounds(#1)}}

\newcommand{\interval}[2]{#1..#2}
\newcommand{\tuple}[1]{\langle #1 \rangle}
\newcommand{\sequence}[1]{[#1]}
\newcommand{\set}[1]{\{#1\}}
\newcommand{\setComp}[2]{\set{#1 ~|~ #2}}
\newcommand{\sequenceComp}[2]{\sequence{#1 ~|~ #2}}
\newcommand{\dom}[1]{D(#1)}
\newcommand{\pp}{+\!\!+}

\newcommand{\rhs}{s}

% ── Table-specific notation ──────────────────────────────────────────────────
\newcommand{\cols}{n}
\newcommand{\rows}{m}
\newcommand{\row}{r}
\newcommand{\col}{j}

\newcommand{\llinex}[1]{Line \ref{#1}}
\newcommand{\linex}[1]{line \ref{#1}}
\newcommand{\lines}[2]{lines \ref{#1}--\ref{#2}}

\newcommand{\ttable}{\texttt{table}\xspace}
\newcommand{\btable}{T^{\Bools}}
\newcommand{\btableTransposed}{\hat{T}^{\Bools}}
\newcommand{\transpose}[1]{\hat{#1}}

\newcommand{\nolbbd}[2]{#2}

% ── Configuration / approach labels ──────────────────────────────────────────
\newcommand{\textconfig}[1]{\textsc{#1}}
\newcommand{\baseGleb}{\textconfig{Int}\xspace}
\newcommand{\baseBool}{\textconfig{Bool}\xspace}

\newcommand{\baseMdd}{\textconfig{MDD}\xspace}
\newcommand{\baseMddInput}{\textconfig{\baseMdd-input}\xspace}
\newcommand{\baseMddDomIncr}{\textconfig{\baseMdd-dom}\xspace}
\newcommand{\baseMddGreedy}{\textconfig{\baseMdd-greedy}\xspace}

\newcommand{\lazyGenerate}{\textconfig{Generate}\xspace}
\newcommand{\lazyNone}{\textconfig{Generate}\xspace}
\newcommand{\lazyFractional}{\textconfig{Fractional}\xspace}
\newcommand{\lazyNegatives}{\textconfig{Negatives}\xspace}
\newcommand{\lazyCutlift}{\textconfig{Cutlift}\xspace}
\newcommand{\lazyShrink}{\textconfig{Shrink}\xspace}
\newcommand{\lazyCutoff}{\textconfig{Hybrid}\xspace}
\newcommand{\lazyCutoffT}[1]{\textconfig{Hybrid (#1)}\xspace}
\newcommand{\lazyAll}{\textconfig{All}\xspace}

% ── xy-pic node macros ───────────────────────────────────────────────────────
\newcommand{\xyo}[1]{*++[o][F-]{#1}}
\newcommand{\xyob}[1]{*++[o][F=]{#1}}
\newcommand{\xyoy}[1]{*++[o][F=]{#1}}
\newcommand{\xyoo}[1]{*++[o][F=]{#1}}

% ── MDD node labels (shared by \mddDiagram*, \mddFlowA--D) ────────────────────
\newcommand{\mddNodeRoot}{[]}
\newcommand{\mddNodeXtwo}{[2]}
\newcommand{\mddNodeXone}{[1]}
\newcommand{\mddNodeXfour}{[4]}
\newcommand{\mddNodeXthree}{[3]}
\newcommand{\mddNodeXtwoYone}{[2,1]}
\newcommand{\mddNodeMerge}{\scriptscriptstyle[1,2] \scriptscriptstyle[4,3]}
\newcommand{\mddNodeXthreeYtwo}{[3,2]}
\newcommand{\mddNodeSnk}{snk}

% ── Algorithm macros ─────────────────────────────────────────────────────────
\newcommand{\codeBlockDelim}{}
\newcommand{\tableCon}{\ensuremath{\texttt{table}(\hat{b}, \btable)}}
\newcommand{\generateInitArgs}{\ensuremath{\assignment, \tableCon}}
\newcommand{\generateArgs}{\ensuremath{\generateInitArgs, X, \rhs, R, V}}
\newcommand{\validCut}{\sum_{\col \in X} a_\col b_\col \leq \rhs}

% ── Cut lifting macros ──────────────────────────────────────────────────────
\newcommand{\Cols}{\interval{1}{\cols}}
\newcommand{\Rows}{\interval{1}{\rows}}
\newcommand{\colLift}{\col'}
\newcommand{\Rtight}{\ensuremath{R_{\text{tight}}}}
\newcommand{\Ntight}{\ensuremath{N_{\text{tight}}}}
\newcommand{\Xtight}{\ensuremath{X'}}
\newcommand{\cutliftThm}{%
Lifting is correct.
Suppose $\btable \models \sum_{\col \in X} a_\col b_\col \leq \rhs$  and $\btable \models b_{\colLift} \rightarrow \sum_{\col \in X} a_\col b_\col \leq \rhs-a_{\colLift}$ then
$\btable \models \sum_{\col \in X \cup \set{\colLift}} a_\col b_\col \leq \rhs$.%
}

% ── Negatives macros ─────────────────────────────────────────────────────────
\newcommand{\negate}[1]{\overline{#1}}
\newcommand{\partSize}[1]{\hat{P}_{#1}}
\newcommand{\unencoded}{B}

% ── Experimental evaluation ──────────────────────────────────────────────────
\newcommand{\timeout}{600\xspace}
\newcommand{\memout}{8\xspace}
\newcommand{\unk}{\texttt{UNK}\xspace}
\newcommand{\mem}{\texttt{MEM}\xspace}
\newcommand{\pst}{\texttt{POST}\xspace}
\newcommand{\sat}{\texttt{FEAS}\xspace}
\newcommand{\uns}{\texttt{UNS}\xspace}
\newcommand{\sol}{\texttt{SOLVED}\xspace}
\newcommand{\opt}{\texttt{OPT}\xspace}
\newcommand{\XCSPYears}{XCSP3'22\textendash25\xspace}
\newcommand{\thousands}{$10^3$}

% ── Figures ───────────────────────────────────────────────────────────────────
% Locally thicken xy-pic edges/frames within a group (restored at the closing brace)
\makeatletter
\newcommand{\xyThick}{\begingroup\xylinethick@=\dimexpr2\xylinethick@\relax}
\newcommand{\xyThickEnd}{\endgroup}
\makeatother
% MDD diagram for the running example (parameterised by xy spacing)
\newcommand{\mddDiagram}[2]{% #1 = @C (column sep), #2 = @R (row sep)
\xyThick
\xymatrix@C=#1@R=#2{
   \xyo{\mddNodeRoot}
       \ar[r]^{\be{1}{2}}
       \ar[rd]^{\be{1}{1}}
       \ar[rdd]^>>>>>>{\be{1}{4}}
       \ar[rddd]_{\be{1}{3}}
       & \xyo{\mddNodeXtwo} \ar[r]^{\be{2}{1}}
       & \xyo{\mddNodeXtwoYone} \ar[r]^{\be{3}{1}} \ar@/_2pc/[r]^{\be{3}{2}}
       & \xyoo{\mddNodeSnk} \\
   & \xyo{\mddNodeXone} \ar[rd]^{\be{2}{2}} &  \\
   & \xyo{\mddNodeXfour} \ar[r]_{\be{2}{3}}
       & \xyo{\mddNodeMerge} \ar[ruu]^{\be{3}{3}} \\
   & \xyo{\mddNodeXthree} \ar[r]^{\be{2}{2}}
       & \xyo{\mddNodeXthreeYtwo} \ar[ruuu]_{\be{3}{2}} \\
}%
\xyThickEnd
}
\newcommand{\mddDiagramHighlight}[2]{% same but with path 1,2,3 highlighted
\xyThick
\xymatrix@C=#1@R=#2{
   \xyo{\mddNodeRoot}
       \ar[r]^{\be{1}{2}}
       \ar@[blue][rd]^{\textcolor{blue}{\be{1}{1}}}
       \ar[rdd]^>>>>>>{\be{1}{4}}
       \ar[rddd]_{\be{1}{3}}
       & \xyo{\mddNodeXtwo} \ar[r]^{\be{2}{1}}
       & \xyo{\mddNodeXtwoYone} \ar[r]^{\be{3}{1}} \ar@/_2pc/[r]^{\be{3}{2}}
       & \xyoo{\mddNodeSnk} \\
   & \xyo{\textcolor{blue}{\mddNodeXone}} \ar@[blue][rd]^{\textcolor{blue}{\be{2}{2}}} &  \\
   & \xyo{\mddNodeXfour} \ar[r]_{\be{2}{3}}
       & \xyo{\textcolor{blue}{\mddNodeMerge}} \ar@[blue][ruu]^{\textcolor{blue}{\be{3}{3}}} \\
   & \xyo{\mddNodeXthree} \ar[r]^{\be{2}{2}}
       & \xyo{\mddNodeXthreeYtwo} \ar[ruuu]_{\be{3}{2}} \\
}%
\xyThickEnd
}
\newcommand{\mddDiagramHighlightBoth}[2]{% paths 1,2,3 (blue) and 4,3,3 (red)
\xyThick
\xymatrix@C=#1@R=#2{
   \xyo{\mddNodeRoot}
       \ar[r]^{\be{1}{2}}
       \ar@[blue][rd]^{\textcolor{blue}{\be{1}{1}}}
       \ar@[red][rdd]^>>>>>>{\textcolor{red}{\be{1}{4}}}
       \ar[rddd]_{\be{1}{3}}
       & \xyo{\mddNodeXtwo} \ar[r]^{\be{2}{1}}
       & \xyo{\mddNodeXtwoYone} \ar[r]^{\be{3}{1}} \ar@/_2pc/[r]^{\be{3}{2}}
       & \xyoo{\mddNodeSnk} \\
   & \xyo{\textcolor{blue}{\mddNodeXone}} \ar@[blue][rd]^{\textcolor{blue}{\be{2}{2}}} &  \\
   & \xyo{\textcolor{red}{\mddNodeXfour}} \ar@[red][r]_{\textcolor{red}{\be{2}{3}}}
       & \xyo{\textcolor{violet}{\mddNodeMerge}} \ar@[violet][ruu]^{\textcolor{violet}{\be{3}{3}}} \\
   & \xyo{\mddNodeXthree} \ar[r]^{\be{2}{2}}
       & \xyo{\mddNodeXthreeYtwo} \ar[ruuu]_{\be{3}{2}} \\
}%
\xyThickEnd
}

% ── Misc ─────────────────────────────────────────────────────────────────────
\newcommand{\omitVarJ}{[x_1, \ldots, x_{j-1}, x_{j+1}, \ldots, x_{n-1}, x_n]}
 after \documentclass and package loading

% ── Glossary / acronyms ─────────────────────────────────────────────────────
\newacronym{sat}{SAT}{\emph{Boolean Satisfiability}}
\newacronym{maxSat}{maxSAT}{\emph{Maximum Satisfiability}}
\newacronym{pb}{PB}{\emph{Pseudo-Boolean}}
\newacronym{cp}{CP}{\emph{Constraint Programming}}
\newacronym{tsp}{TSP}{\emph{Travelling Salesperson Problem}}
\newacronym{mdd}{MDD}{\emph{Multivalued Decision Diagram}}
\newacronym{bdd}{BDD}{\emph{Binary Decision Diagram}}
\newacronym{csp}{CSP}{\emph{Constraint Satisfaction Problem}}
\newacronym{cop}{COP}{\emph{Constrained Optimisation Problem}}
\newacronym{mip}{MIP}{\emph{Mixed Integer Programming}}
\newacronym{ilp}{ILP}{\emph{Integer Linear Programming}}
\newacronym{lp}{LP}{\emph{Linear Programming}}
\newacronym{lbbd}{LBBD}{\emph{Logic-Based Benders Decomposition}}
\newacronym{mus}{MUS}{\emph{Minimal Unsatisfiable Subset}}

\newcommand{\milp}{\gls{ilp}\xspace}
\newcommand{\relaxation}{\acrshort{lp} relaxation\xspace}

% ── Text shorthands ──────────────────────────────────────────────────────────
\newcommand{\dquote}[1]{``#1''}
\newcommand{\quoted}[1]{`#1'}
\newcommand{\ie}{i.e.\xspace}
\newcommand{\eg}{e.g.\xspace}
\newcommand{\Eg}{E.g.\xspace}

% ── Author annotation commands (silenced by default) ─────────────────────────
\newcommand{\question}[1]{}
\newcommand{\pjs}[1]{}
\newcommand{\tias}[1]{}
\newcommand{\hb}[1]{}
\newcommand{\wout}[1]{}
\newcommand{\scrap}[1]{}
\newcommand{\ignore}[1]{}

\newcommand{\red}[1]{\textcolor{blue}{#1}}

% ── Math notation ────────────────────────────────────────────────────────────
\newcommand{\brackets}[1]{\ensuremath{[\![#1]\!]}}
\newcommand{\iverson}[1]{\ensuremath{\langle #1 \rangle}}
\newcommand{\beq}[2]{\brackets{#1 = #2}}
\newcommand{\bevarlookup}[1]{\ifnum9<1#1 \ifcase#1\or x\or y\or z\fi\else x_{#1}\fi}
\newcommand{\x}[1]{\bevarlookup{#1}}
\newcommand{\be}[2]{\beq{\bevarlookup{#1}}{#2}}
\newcommand{\besml}[2]{\scriptstyle\be{#1}{#2}}

\newcommand{\etAl}[1]{#1~\emph{et al.}}
\newcommand{\citeNoun}[2]{\etAl{#1}~\cite{#2}}

\DeclareMathOperator*{\argmax}{arg\,max}
\DeclareMathOperator*{\argmin}{arg\,min}

\newcommand{\assignment}{\ensuremath{\hat{v}}}
\newcommand{\true}{\mathit{true}}
\newcommand{\false}{\mathit{false}}
\newcommand{\Integers}{\mathbb{Z}}
\newcommand{\Reals}{\mathbb{R}}
\newcommand{\Bools}{\mathbb{B}}
\renewcommand{\emptyset}{\varnothing}

\newcommand{\vars}[1]{\ensuremath{vars(#1)}}
\newcommand{\scope}[1]{\ensuremath{scope(#1)}}
\newcommand{\cons}[1]{\ensuremath{cons(#1)}}
\newcommand{\bounds}[1]{\ensuremath{bounds(#1)}}

\newcommand{\interval}[2]{#1..#2}
\newcommand{\tuple}[1]{\langle #1 \rangle}
\newcommand{\sequence}[1]{[#1]}
\newcommand{\set}[1]{\{#1\}}
\newcommand{\setComp}[2]{\set{#1 ~|~ #2}}
\newcommand{\sequenceComp}[2]{\sequence{#1 ~|~ #2}}
\newcommand{\dom}[1]{D(#1)}
\newcommand{\pp}{+\!\!+}

\newcommand{\rhs}{s}

% ── Table-specific notation ──────────────────────────────────────────────────
\newcommand{\cols}{n}
\newcommand{\rows}{m}
\newcommand{\row}{r}
\newcommand{\col}{j}

\newcommand{\llinex}[1]{Line \ref{#1}}
\newcommand{\linex}[1]{line \ref{#1}}
\newcommand{\lines}[2]{lines \ref{#1}--\ref{#2}}

\newcommand{\ttable}{\texttt{table}\xspace}
\newcommand{\btable}{T^{\Bools}}
\newcommand{\btableTransposed}{\hat{T}^{\Bools}}
\newcommand{\transpose}[1]{\hat{#1}}

\newcommand{\nolbbd}[2]{#2}

% ── Configuration / approach labels ──────────────────────────────────────────
\newcommand{\textconfig}[1]{\textsc{#1}}
\newcommand{\baseGleb}{\textconfig{Int}\xspace}
\newcommand{\baseBool}{\textconfig{Bool}\xspace}

\newcommand{\baseMdd}{\textconfig{MDD}\xspace}
\newcommand{\baseMddInput}{\textconfig{\baseMdd-input}\xspace}
\newcommand{\baseMddDomIncr}{\textconfig{\baseMdd-dom}\xspace}
\newcommand{\baseMddGreedy}{\textconfig{\baseMdd-greedy}\xspace}

\newcommand{\lazyGenerate}{\textconfig{Generate}\xspace}
\newcommand{\lazyNone}{\textconfig{Generate}\xspace}
\newcommand{\lazyFractional}{\textconfig{Fractional}\xspace}
\newcommand{\lazyNegatives}{\textconfig{Negatives}\xspace}
\newcommand{\lazyCutlift}{\textconfig{Cutlift}\xspace}
\newcommand{\lazyShrink}{\textconfig{Shrink}\xspace}
\newcommand{\lazyCutoff}{\textconfig{Hybrid}\xspace}
\newcommand{\lazyCutoffT}[1]{\textconfig{Hybrid (#1)}\xspace}
\newcommand{\lazyAll}{\textconfig{All}\xspace}

% ── xy-pic node macros ───────────────────────────────────────────────────────
\newcommand{\xyo}[1]{*++[o][F-]{#1}}
\newcommand{\xyob}[1]{*++[o][F=]{#1}}
\newcommand{\xyoy}[1]{*++[o][F=]{#1}}
\newcommand{\xyoo}[1]{*++[o][F=]{#1}}

% ── MDD node labels (shared by \mddDiagram*, \mddFlowA--D) ────────────────────
\newcommand{\mddNodeRoot}{[]}
\newcommand{\mddNodeXtwo}{[2]}
\newcommand{\mddNodeXone}{[1]}
\newcommand{\mddNodeXfour}{[4]}
\newcommand{\mddNodeXthree}{[3]}
\newcommand{\mddNodeXtwoYone}{[2,1]}
\newcommand{\mddNodeMerge}{\scriptscriptstyle[1,2] \scriptscriptstyle[4,3]}
\newcommand{\mddNodeXthreeYtwo}{[3,2]}
\newcommand{\mddNodeSnk}{snk}

% ── Algorithm macros ─────────────────────────────────────────────────────────
\newcommand{\codeBlockDelim}{}
\newcommand{\tableCon}{\ensuremath{\texttt{table}(\hat{b}, \btable)}}
\newcommand{\generateInitArgs}{\ensuremath{\assignment, \tableCon}}
\newcommand{\generateArgs}{\ensuremath{\generateInitArgs, X, \rhs, R, V}}
\newcommand{\validCut}{\sum_{\col \in X} a_\col b_\col \leq \rhs}

% ── Cut lifting macros ──────────────────────────────────────────────────────
\newcommand{\Cols}{\interval{1}{\cols}}
\newcommand{\Rows}{\interval{1}{\rows}}
\newcommand{\colLift}{\col'}
\newcommand{\Rtight}{\ensuremath{R_{\text{tight}}}}
\newcommand{\Ntight}{\ensuremath{N_{\text{tight}}}}
\newcommand{\Xtight}{\ensuremath{X'}}
\newcommand{\cutliftThm}{%
Lifting is correct.
Suppose $\btable \models \sum_{\col \in X} a_\col b_\col \leq \rhs$  and $\btable \models b_{\colLift} \rightarrow \sum_{\col \in X} a_\col b_\col \leq \rhs-a_{\colLift}$ then
$\btable \models \sum_{\col \in X \cup \set{\colLift}} a_\col b_\col \leq \rhs$.%
}

% ── Negatives macros ─────────────────────────────────────────────────────────
\newcommand{\negate}[1]{\overline{#1}}
\newcommand{\partSize}[1]{\hat{P}_{#1}}
\newcommand{\unencoded}{B}

% ── Experimental evaluation ──────────────────────────────────────────────────
\newcommand{\timeout}{600\xspace}
\newcommand{\memout}{8\xspace}
\newcommand{\unk}{\texttt{UNK}\xspace}
\newcommand{\mem}{\texttt{MEM}\xspace}
\newcommand{\pst}{\texttt{POST}\xspace}
\newcommand{\sat}{\texttt{FEAS}\xspace}
\newcommand{\uns}{\texttt{UNS}\xspace}
\newcommand{\sol}{\texttt{SOLVED}\xspace}
\newcommand{\opt}{\texttt{OPT}\xspace}
\newcommand{\XCSPYears}{XCSP3'22\textendash25\xspace}
\newcommand{\thousands}{$10^3$}

% ── Figures ───────────────────────────────────────────────────────────────────
% Locally thicken xy-pic edges/frames within a group (restored at the closing brace)
\makeatletter
\newcommand{\xyThick}{\begingroup\xylinethick@=\dimexpr2\xylinethick@\relax}
\newcommand{\xyThickEnd}{\endgroup}
\makeatother
% MDD diagram for the running example (parameterised by xy spacing)
\newcommand{\mddDiagram}[2]{% #1 = @C (column sep), #2 = @R (row sep)
\xyThick
\xymatrix@C=#1@R=#2{
   \xyo{\mddNodeRoot}
       \ar[r]^{\be{1}{2}}
       \ar[rd]^{\be{1}{1}}
       \ar[rdd]^>>>>>>{\be{1}{4}}
       \ar[rddd]_{\be{1}{3}}
       & \xyo{\mddNodeXtwo} \ar[r]^{\be{2}{1}}
       & \xyo{\mddNodeXtwoYone} \ar[r]^{\be{3}{1}} \ar@/_2pc/[r]^{\be{3}{2}}
       & \xyoo{\mddNodeSnk} \\
   & \xyo{\mddNodeXone} \ar[rd]^{\be{2}{2}} &  \\
   & \xyo{\mddNodeXfour} \ar[r]_{\be{2}{3}}
       & \xyo{\mddNodeMerge} \ar[ruu]^{\be{3}{3}} \\
   & \xyo{\mddNodeXthree} \ar[r]^{\be{2}{2}}
       & \xyo{\mddNodeXthreeYtwo} \ar[ruuu]_{\be{3}{2}} \\
}%
\xyThickEnd
}
\newcommand{\mddDiagramHighlight}[2]{% same but with path 1,2,3 highlighted
\xyThick
\xymatrix@C=#1@R=#2{
   \xyo{\mddNodeRoot}
       \ar[r]^{\be{1}{2}}
       \ar@[blue][rd]^{\textcolor{blue}{\be{1}{1}}}
       \ar[rdd]^>>>>>>{\be{1}{4}}
       \ar[rddd]_{\be{1}{3}}
       & \xyo{\mddNodeXtwo} \ar[r]^{\be{2}{1}}
       & \xyo{\mddNodeXtwoYone} \ar[r]^{\be{3}{1}} \ar@/_2pc/[r]^{\be{3}{2}}
       & \xyoo{\mddNodeSnk} \\
   & \xyo{\textcolor{blue}{\mddNodeXone}} \ar@[blue][rd]^{\textcolor{blue}{\be{2}{2}}} &  \\
   & \xyo{\mddNodeXfour} \ar[r]_{\be{2}{3}}
       & \xyo{\textcolor{blue}{\mddNodeMerge}} \ar@[blue][ruu]^{\textcolor{blue}{\be{3}{3}}} \\
   & \xyo{\mddNodeXthree} \ar[r]^{\be{2}{2}}
       & \xyo{\mddNodeXthreeYtwo} \ar[ruuu]_{\be{3}{2}} \\
}%
\xyThickEnd
}
\newcommand{\mddDiagramHighlightBoth}[2]{% paths 1,2,3 (blue) and 4,3,3 (red)
\xyThick
\xymatrix@C=#1@R=#2{
   \xyo{\mddNodeRoot}
       \ar[r]^{\be{1}{2}}
       \ar@[blue][rd]^{\textcolor{blue}{\be{1}{1}}}
       \ar@[red][rdd]^>>>>>>{\textcolor{red}{\be{1}{4}}}
       \ar[rddd]_{\be{1}{3}}
       & \xyo{\mddNodeXtwo} \ar[r]^{\be{2}{1}}
       & \xyo{\mddNodeXtwoYone} \ar[r]^{\be{3}{1}} \ar@/_2pc/[r]^{\be{3}{2}}
       & \xyoo{\mddNodeSnk} \\
   & \xyo{\textcolor{blue}{\mddNodeXone}} \ar@[blue][rd]^{\textcolor{blue}{\be{2}{2}}} &  \\
   & \xyo{\textcolor{red}{\mddNodeXfour}} \ar@[red][r]_{\textcolor{red}{\be{2}{3}}}
       & \xyo{\textcolor{violet}{\mddNodeMerge}} \ar@[violet][ruu]^{\textcolor{violet}{\be{3}{3}}} \\
   & \xyo{\mddNodeXthree} \ar[r]^{\be{2}{2}}
       & \xyo{\mddNodeXthreeYtwo} \ar[ruuu]_{\be{3}{2}} \\
}%
\xyThickEnd
}

% ── Misc ─────────────────────────────────────────────────────────────────────
\newcommand{\omitVarJ}{[x_1, \ldots, x_{j-1}, x_{j+1}, \ldots, x_{n-1}, x_n]}


\newcommand{\compare}{base_gurobi-mdd-reduce-dom-incr-nocombined--lazy_gurobi-hybrid_0}


\title{Table Constraints for Integer Programming (CP'26)}
\subtitle{JFPC'26}
\author[\etAl{H. Bierlee}]{Hendrik Bierlee\inst{1} \and Wout Piessens\inst{1} \and Tias Guns\inst{1} \and Peter J. Stuckey\inst{2,3}}
\institute[KU Leuven, Monash University]{\inst{1} KU Leuven, Belgium \and \inst{2} Monash University, Melbourne, Australia \and \inst{3} OPTIMA ITTC, Melbourne, Australia}
\date{}

\titlegraphic{%
\includegraphics[height=1.2cm]{erc-logo.pdf}\hspace{1cm}%
\includegraphics[height=1.2cm]{res/kuleuven.png}%
}

\begin{document}

\begin{frame}
\titlepage
\begin{textblock*}{12cm}(2cm,8.25cm)
\centering{\scriptsize Funded by ERC (CHAT-Opt, 101002802), FWO (G020525N), and ARC (OPTIMA).}
\end{textblock*}
\begin{textblock*}{2cm}(0.3cm,5.5cm)
\centering\includesvg[width=1.5cm]{qr-paper}\\[-4pt]{\tiny Paper}
\end{textblock*}
\begin{textblock*}{2cm}(13.8cm,5.5cm)
\centering\includesvg[width=1.5cm]{qr-slides}\\[-4pt]{\tiny Slides}
\end{textblock*}
\end{frame}

% \begin{frame}{Outline}
% \tableofcontents
% \end{frame}

%===============================================================================
\section{Introduction}
%===============================================================================

\begin{frame}{Table Constraints in \gls{cp}}

\begin{itemize}
	\item Solver-independent modelling languages (\eg MiniZinc~\cite{minizinc}, CPMpy~\cite{guns2019increasing}) solve \gls{cp} models using \gls{ilp} solvers
    \item But there is a large perfomance gap between \gls{ilp} and \eg CP-SAT
    \begin{itemize}
	    \item Because globals (\eg table constraints) are not linearized as well
	    \item Because we do not use the full \gls{ilp} toolset
    \end{itemize}
\end{itemize}

% \begin{itemize}
%     \item Table constraints \texttt{table}($\hat{x}, T$) enforces that variables $\hat{x}$ take values from a row in $T$
%     \item One of the most well-studied global constraints in \gls{cp}~\cite{GAC,STR2}
%     \item Prevalent in configuration, scheduling, in $\approx 25\%$ of all XCSP3 competition instances~\cite{audemard2025proceedings2025xcsp3competition}
% \end{itemize}

\medskip
	\textbf{Running example:} \texttt{table}($[\x{1}, \x{2}, \x{3}], T$) with enforces $\hat{x}$ to take one row of $T$:
\[
\begin{array}{rrr}
\x{1} & \x{2} & \x{3} \\ \hline
2 & 1 & 1 \\
3 & 2 & 2 \\
4 & 3 & 3 \\
1 & 2 & 3 \\
2 & 1 & 2
\end{array}
\]
\end{frame}

%===============================================================================
\section{Eager Encodings}
%===============================================================================

\begin{frame}{Simple Integer Encoding (\baseGleb)}
\begin{itemize}
    \item Standard encoding by \citeNoun{Belov}{belov_improved_2016} (default in MiniZinc)
    \item Binary variables $r_i \in \set{0,1}$ for each row $i$, s.t. $r_i = 1$ iff row $i$ chosen
    % \item Compact: $m$ new variables, $k{+}1$ constraints (with $m$ number of rows, $k$ number of columns)
\end{itemize}

\begin{columns}[T]
\begin{column}{0.18\textwidth}
\textbf{Table $T$:}\\[0.5em]
\begin{tabular}{ccc}
$\x{1}$ & $\x{2}$ & $\x{3}$ \\ \hline
2 & 1 & 1 \\
3 & 2 & 2 \\
4 & 3 & 3 \\
1 & 2 & 3 \\
2 & 1 & 2
\end{tabular}
\end{column}
\begin{column}{0.37\textwidth}
\textbf{Encoding:}
\begin{align*}
\sum_{l \in 1..m} r_l  &= 1 \\
\sum_{l \in 1..m} T_{l,i} r_l &= x_i, \quad i \in 1..k
\end{align*}
\end{column}
\begin{column}{0.43\textwidth}
\textbf{Example:}
\[
\begin{array}{rrrrrcl}
 r_1 &+  r_2& +  r_3& + r_4& + r_5 &= &1 \\
2 r_1& + 3r_2& + 4 r_3& + 1 r_4& + 2r_5 &=& \x{1} \\
1 r_1& + 2r_2& + 3 r_3& + 2r_4&+ 1 r_5 &=& \x{2} \\
1 r_1& + 2r_2& + 3 r_3& + 3r_4&+ 2r_5 &=& \x{3}
\end{array}
\]
\end{column}
\end{columns}

\vfill
	For example, if we choose row $1$, then $r_1 = 1$, so $\x{1} = 2 \cdot r_1 + 0 + \ldots + 0 = 2$

\end{frame}


\begin{frame}{Booleanized Encoding (\baseBool)~\cite{petit}}

	We booleanize $\x{1}$ with binary variables $\be{1}{j} \in \set{0,1}, j \in D(\x{1})$ such that $\be{1}{j} \Leftrightarrow \x{1} = j$:

	\[
	\sum_{j \in D(\x{1})} \be{1}{j} = 1, \qquad
	\sum_{j \in D(\x{1})} j \cdot \be{1}{j} = \x{1}
	\]


% \begin{itemize}
% 	\item The row index variable is \textit{nominal}, and so are the column variables $x_i$
%     % \item Key insight: table variables are often \emph{nominal} -- other constraints reason about $x_i = v$
%     \item Direct (unary) encoding: binary variable $\be{i}{j}$ for each value $j \in \dom{x_i}$
%     \item Rewrite integer table into Boolean 01 table $\btable$
% \end{itemize}

	\vfill
	And booleanize $T$ as a binary table $\btable$:
	% the table $\texttt{table}(\sequence{\be{1}{1},\be{1}{2}, \ldots, \be{3}{3}}, \btable)$ as well:

\[
\begin{array}{rrr}
\x{1} & \x{2} & \x{3} \\ \hline
2 & 1 & 1 \\
3 & 2 & 2 \\
4 & 3 & 3 \\
1 & 2 & 3 \\
2 & 1 & 2
\end{array}
\quad \equiv \quad
\begin{array}{l|rrrr|rrr|rrr}
&  {\besml{1}{1}} & \besml{1}{2} & \besml{1}{3} & \besml{1}{4} & \besml{2}{1} & \besml{2}{2} & \besml{2}{3} & \besml{3}{1} & \besml{3}{2} & \besml{3}{3} \\ \hline
&   0   & 1   & 0   & 0   & 1   & 0   & 0   & 1   & 0   & 0 \\
&   0   & 0   & 1   & 0   & 0   & 1   & 0   & 0   & 1   & 0 \\
&   0   & 0   & 0   & 1   & 0   & 0   & 1   & 0   & 0   & 1 \\
&   1   & 0   & 0   & 0   & 0   & 1   & 0   & 0   & 0   & 1 \\
&   0   & 1   & 0   & 0   & 1   & 0   & 0   & 0   & 1   & 0
\end{array}
\]

% \smallskip
% Equivalent \gls{lp} relaxation to \baseGleb, but \textbf{stronger model} when binary variables used elsewhere.
\end{frame}

\begin{frame}{MDD-based Flow Encoding (\baseMdd)}

We can reformulate $\btable$ as a reduced Multi-valued Decision Diagram (MDD)~\cite{DBLP:journals/constraints/ChengY10}

\[
\vcenter{\hbox{$\begin{array}{rrr}
\x{1} & \x{2} & \x{3} \\ \hline
2 & 1 & 1 \\
3 & 2 & 2 \\
\alt<3->{\textcolor{red}{4} & \textcolor{red}{3} & \textcolor{red}{3}}{4 & 3 & 3} \\
\alt<2->{\textcolor{blue}{1} & \textcolor{blue}{2} & \textcolor{blue}{3}}{1 & 2 & 3} \\
2 & 1 & 2
\end{array}$}}
\;\equiv\;
\vcenter{\hbox{$\btable$}}
\;\equiv\;
\vcenter{\hbox{$\alt<3->{\mddDiagramHighlightBoth{12mm}{8mm}}{\alt<2>{\mddDiagramHighlight{12mm}{8mm}}{\mddDiagram{12mm}{8mm}}}$}}
\;\equiv\;
	\text{flow constraints}
% \vcenter{\hbox{$\scriptsize\begin{array}{rcl@{~~}l}
% \be{1}{2} &= &\be{2}{1} & ([2])\\
% \be{2}{1} &= &\be{3}{1} + e_{([2,1],2)} &([2,1]) \\
% \be{1}{1} &= &e_{([1],2)} & ([1])\\
% e_{([1],2)} + \be{2}{3} &=& \be{3}{3} & ([1,2], [4,3]) \\
% \be{1}{4} & = & \be{2}{3} & ([4]) \\
% \be{1}{3} &=& e_{([3],2)}  & ([3])\\
% e_{([3],2)} &=& e_{([3,2],2)} & ([3,2])\\
% \be{1}{1} + \be{1}{2} + \be{1}{3} + \be{1}{4} &= &1 & ([]) \\
% e_{([2,1],2)} + \be{3}{1} + \be{3}{3} + e_{([3,2],2)} &=& 1 & (snk)\\
% e_{([2,1],2)} + e_{([3,2],2)} & = &\be{3}{2} \\
% e_{([1],2)} + e_{([3],2)} & = &\be{2}{2} \\
% \end{array}$}}
\]

	The MDD can be \textbf{exponentially smaller}, the flow constraints are \textbf{totally unimodular}

\ignore{


\begin{columns}[T]
\begin{column}{0.45\textwidth}



% \begin{itemize}
%     \item Encode table as a \textbf{reduced \gls{mdd}}~\cite{DBLP:journals/constraints/ChengY10} with nodes sharing prefixes
%     \item Encode (possibly exponentially smaller) \gls{mdd} as a (totally unimodular~\cite{ahuja1993network}) \textbf{flow network} (totally unimodular)
% \end{itemize}

\medskip
Column ordering of the table matters:
\begin{itemize}
    \item \baseMddInput: original order
    \item \baseMddDomIncr: increasing domain size
\end{itemize}
\end{column}
\begin{column}{0.5\textwidth}
\end{columns}

	}
\end{frame}

\ignore{

\begin{frame}{\gls{mdd} Flow Encoding -- Example}
\begin{columns}[T]
\begin{column}{0.45\textwidth}
\centering
\small
\[
\mddDiagram{8mm}{7mm}
\]
\end{column}
\begin{column}{0.52\textwidth}
\textbf{Flow constraints:}
\smallskip
\small
\[
\begin{array}{rcl@{~~}l}
\be{1}{2} &= &\be{2}{1} & \text{node } [2]\\
\be{2}{1} &= &\be{3}{1} + e_{([2,1],2)} & \text{node } [2,1] \\
\be{1}{1} &= &e_{([1],2)} & \text{node } [1]\\
e_{([1],2)} + \be{2}{3} &=& \be{3}{3} & \text{node } [1,2]\!=\![4,3] \\
\be{1}{4} & = & \be{2}{3} & \text{node } [4] \\
\be{1}{3} &=& e_{([3],2)}  & \text{node } [3]\\
e_{([3],2)} &=& e_{([3,2],2)} & \text{node } [3,2]\\[3pt]
\multicolumn{3}{l}{\be{1}{1} + \be{1}{2} + \be{1}{3} + \be{1}{4} = 1} & \text{source}\\
\multicolumn{3}{l}{e_{([2,1],2)} + \be{3}{1} + \be{3}{3} + e_{([3,2],2)} = 1} & \text{sink}\\[3pt]
e_{([2,1],2)} + e_{([3,2],2)} & = &\be{3}{2} & \text{channel}\\
e_{([1],2)} + e_{([3],2)} & = &\be{2}{2} & \text{channel}
\end{array}
\]
\end{column}
\end{columns}

\smallskip
\begin{itemize}
    \item Each \gls{mdd} node $\rightarrow$ flow balance constraint (flow in = flow out)
    \item Source/sink: total flow = 1 (exactly one row active)
    \item Channel: edge variables $\leftrightarrow$ Boolean encoding variables $\be{i}{j}$
\end{itemize}
\end{frame}


\begin{frame}{Encoding Comparison -- Theory}
\begin{center}
\begin{tabular}{lccc}
\toprule
 & \baseGleb & \baseBool & \baseMdd \\
\midrule
New variables & $m$ & $m + \sum |D(x_i)|$ & $|EV|$ \\
Constraints & $k + 1$ & $\sum |D(x_i)| + 1$ & $|VM| + \sum |D(x_i)|$ \\
\gls{lp} relaxation & \multicolumn{2}{c}{equivalent} & equivalent to \baseBool \\
Totally unimodular & no & no & \textbf{yes} (flow) \\
Size vs.\ table & $O(mk)$ & $O(mk)$ & up to \textbf{exp.\ smaller} \\
\bottomrule
\end{tabular}
\end{center}

\bigskip
All three encodings give the same \gls{lp} relaxation on $\hat{x}$, but:
\begin{itemize}
    \item \baseBool enables stronger reasoning when binary variables appear elsewhere
    \item \baseMdd can be exponentially more compact and is totally unimodular
\end{itemize}
\end{frame}
}

\newcommand{\NA}{--}
\newcommand{\resultheader}{%
\toprule
 & {\pst} & {\sat} & {\sol} & {time} & {constraints} & {cuts} \\
\midrule}
\newcommand{\resultencoding}[4][]{%
{#1\baseGleb} & {#1 171} & {#1} & {#1 55} & {#1 851.7} & {#1 8249.3} & {#1 \NA} \\
{#2\baseBool} & {#2 147} & {#2} & {#2 69} & {#2 771.5} & {#2 25670.4} & {#2 \NA} \\
{#3\baseMdd} & {#3 154} & {#3} & {#3 67} & {#3 780.8} & {#3 38800.8} & {#3 \NA} \\
{#4\baseMddDomIncr} & {#4 154} & {#4} & {#4 76} & {#4 728.5} & {#4 30884.2} & {#4 \NA} \\}
\newcommand{\resultlazy}[4][]{%
{#1\lazyGenerate} & {#1 174} & {#1} & {#1 36} & {#1 988.6} & {#1 8639.7} & {#1 1857.5} \\
{#2\lazyFractional} & {#2 174} & {#2} & {#2 36} & {#2 980.9} & {#2 8639.7} & {#2 3202.6} \\
{#3\lazyCutlift} & {#3 174} & {#3} & {#3 45} & {#3 925.0} & {#3 8639.7} & {#3 847.1} \\
{#4\lazyAll} & {#4 174} & {#4} & {#4 55} & {#4 871.5} & {#4 8639.7} & {#4 1326.7} \\}
\newcommand{\resulthybrid}[1][]{%
{#1\lazyCutoff (1000)} & {#1 174} & {#1} & {#1 79} & {#1 717.5} & {#1 29635.4} & {#1 49.9} \\
{#1\lazyCutoff (2000)} & {#1 174} & {#1} & {#1\bfseries 80} & {#1\bfseries 715.6} & {#1 29912.8} & {#1 49.9} \\
{#1\lazyCutoff (3000)} & {#1\bfseries 175} & {#1} & {#1 78} & {#1 716.7} & {#1 30243.3} & {#1 37.3} \\}
\newcommand{\resulttabular}[1]{%
\begin{tabular}{lS[table-format=3.0]S[table-format=2.0]S[table-format=2.0]S[table-format=3.1]S[table-format=5.1]S[table-format=4.1]}
\resultheader
#1
\bottomrule
\end{tabular}}
\begin{frame}{Results -- Encodings (CSP, 180 instances)}

\begin{description}
    \item[Benchmarks] XCSP3 competition instances (2022--2025)~\cite{audemard2025proceedings2025xcsp3competition}, 180 CSP instances
    \item[Solver] Gurobi 13.0.1~\cite{gurobi}, single core, 600s time limit, 8GB memory
\end{description}

\vfill

\centering
\small
\resulttabular{%
\resultencoding
{\alt<2->{}{\color{white}}}%
{\alt<3->{}{\color{white}}}%
{\alt<4->{}{\color{white}}}%
}

\smallskip
{\scriptsize pst = posted, sat = feasible, sol = solved. Time = PAR2 avg [s.]}
\end{frame}

%===============================================================================
\section{Lazy Cut Generation}
%===============================================================================

\begin{frame}{Lazy Cut Generation for Tables}

Instead of an eager encoding, we can generated cuts \textbf{lazily} (\ie row generation, logic-based Benders decomposition~\cite{DBLP:journals/ior/Hooker07})

\begin{enumerate}
	\item<1-> Solve without encoding tables, get assignment $\hat{v}$ (\eg $\hat{v} = \sequence{2,2,2}$)
	\item<2-> If $\hat{v}$ violates a table, generate a \textbf{cut} (violating $\hat{v}$ but not the table)
		\begin{itemize}
			\item \Eg $\be{1}{2} + \be{2}{2} + \be{3}{2} \leq 2$
			\item Note this cut is a logical clause $\neg \be{1}{2} \lor \neg \be{2}{2} \lor \neg \be{3}{2}$
		\end{itemize}
    \item<3-> Repeat until $\hat{v}$ satisfies all tables
\end{enumerate}

\onslide<4->{
Can we generate a stronger cut, \eg $\be{1}{2} + \be{2}{2} \leq 1$?
}

%
% \textbf{Example:} Assignment $\hat{v} = (2,2,2)$ is not in the table, so we generate a cut:
% \begin{equation*}
% \be{1}{2} + \be{2}{2} + \be{3}{2} \leq 2
% \end{equation*}
% Note this is a logical clause $\neg \be{1}{2} \lor \neg \be{2}{2} \lor \neg \be{3}{2}$
% % Stronger cut: $\be{1}{2} + \be{2}{2} \leq 1$ \quad (also cuts $(2,2,1)$ and $(2,2,3)$)
\end{frame}

\ignore{

\begin{frame}{The \lazyGenerate Algorithm}
\begin{columns}[T]
\begin{column}{0.55\textwidth}
\begin{itemize}
    \item Start with the \emph{false} cut ($0 \leq -1$)
    \item Iteratively relax by adding one partition (integer variable) at a time
    \item Greedy: choose partition that removes the most rows from $R$
    \[
    \arg\min_{l} |R \cap \bigcup_{\col \in C(\assignment, l)} \btable_\col|
    \]
    \item Stop when no table rows are violated ($R = \emptyset$)
\end{itemize}
\end{column}
\begin{column}{0.42\textwidth}
\textbf{Example:} Cut $\assignment = (2,2,2)$
\begin{enumerate}
    \item $X{=}\emptyset, \rhs{=}{-}1, R{=}\{1..5\}$
    \item Choose $\x{1}$: add $\be{1}{2}$\\
          $\rhs{=}0, R{=}\{1,5\}$
    \item Choose $\x{2}$: add $\be{2}{2}$\\
          $\rhs{=}1, R{=}\emptyset$
    \item Return $\be{1}{2} + \be{2}{2} \leq 1$
\end{enumerate}
\end{column}
\end{columns}
\end{frame}

}

\newcommand{\violated}{\cellcolor{red!20}}
\newcommand{\accepted}{\cellcolor{blue!15}}
\newcommand{\tblcw}{8mm}% shared column width for Boolean table visualizations
\newcommand{\gh}[1]{\alt<2->{\textcolor{black!15}{#1}}{#1}}% ghost: fade from slide 2+
\begin{frame}{\lazyGenerate Example: $\assignment = (2,2,2)$}

\centering
\smallskip
\small
\renewcommand{\arraystretch}{1.15}
\begin{tabular}{c|>{\centering}p{\tblcw}>{\centering}p{\tblcw}>{\centering}p{\tblcw}>{\centering}p{\tblcw}|>{\centering}p{\tblcw}>{\centering}p{\tblcw}>{\centering}p{\tblcw}|>{\centering}p{\tblcw}>{\centering}p{\tblcw}>{\centering\arraybackslash}p{\tblcw}}
 & \gh{$\besml{1}{1}$} & $\besml{1}{2}$ & \gh{$\besml{1}{3}$} & \gh{$\besml{1}{4}$}
 & \gh{$\besml{2}{1}$} & $\besml{2}{2}$ & \gh{$\besml{2}{3}$}
 & \gh{$\besml{3}{1}$} & $\besml{3}{2}$ & \gh{$\besml{3}{3}$} \\
\hline
\only<1>{\violated}\only<2>{\violated}\only<3>{\violated}%
$1$
 & \gh{0} & 1 & \gh{0} & \gh{0}
 & \gh{1} & \only<4->{\accepted}0 & \gh{0}
 & \gh{1} & 0 & \gh{0} \\
\only<1>{\violated}\only<2>{\violated}%
$2$
 & \gh{0} & \only<3->{\accepted}0 & \gh{1} & \gh{0}
 & \gh{0} & 1 & \gh{0}
 & \gh{0} & 1 & \gh{0} \\
\only<1>{\violated}\only<2>{\violated}%
$3$
 & \gh{0} & \only<3->{\accepted}0 & \gh{0} & \gh{1}
 & \gh{0} & \only<4->{\accepted}0 & \gh{1}
 & \gh{0} & 0 & \gh{1} \\
\only<1>{\violated}\only<2>{\violated}%
$4$
 & \gh{1} & \only<3->{\accepted}0 & \gh{0} & \gh{0}
 & \gh{0} & 1 & \gh{0}
 & \gh{0} & 0 & \gh{1} \\
\only<1>{\violated}\only<2>{\violated}\only<3>{\violated}%
$5$
 & \gh{0} & 1 & \gh{0} & \gh{0}
 & \gh{1} & \only<4->{\accepted}0 & \gh{0}
 & \gh{0} & 1 & \gh{0} \\
\hline
\violated $\assignment$
 & \gh{0} & 1 & \gh{0} & \gh{0}
 & \gh{0} & 1 & \gh{0}
 & \gh{0} & 1 & \gh{0} \\
\end{tabular}

\bigskip
\only<1-2>{%
Cut: $0 \leq -1$ \quad {\small (all rows \colorbox{red!20}{violated})}%
}%
\only<3>{%
Choose $\x{1}$, add $\be{1}{2}$: \quad
${\color{blue}\be{1}{2}} \leq 0$ \quad {\small (rows 1,5 still violated)}%
}%
\only<4>{%
Choose $\x{2}$, add $\be{2}{2}$: \quad
${\color{blue}\be{1}{2} + \be{2}{2}} \leq 1$ \quad \checkmark%
}%
\end{frame}

\begin{frame}{Extensions: \lazyCutlift and \lazyFractional}

\begin{description}[labelwidth=0.3cm]
	\item<1->[\lazyCutlift] Strengthen valid cut by \textbf{lifting} variables (\textit{without} incrementing RHS)

\medskip
Input: $\be{1}{2} + \be{2}{2} + \be{3}{2} \leq 2$

\onslide<2->{Output: $\be{1}{1} + \be{1}{2} + 2\be{1}{4} + \be{2}{2} + \be{3}{1} + \be{3}{2} \leq 2$}

\bigskip
\item<3->[\lazyFractional] Generate cuts from \textbf{fractional} LP solutions

\medskip
Input: $\assignment = \sequence{0, 0.5, 0.5, 0, 0, 1, 0, 0.5, 0.5, 0}$

\onslide<4->{Output: $\be{1}{2} + \be{2}{2} \leq 1$}
\end{description}

\end{frame}

\ignore{

\begin{frame}{Extension 2: Cut Lifting (\lazyCutlift)}
\begin{itemize}
    \item After generating a valid cut, \textbf{strengthen} it by adding more terms
    \item Add variable $b_{\col'}$ if it doesn't violate any table row
    \item Key concept: \textbf{tight rows} -- rows where LHS $= \rhs$
    \begin{itemize}
        \item Can only lift variables not appearing in any tight row
        \item New lifts may create new tight rows $\Rightarrow$ iterate
    \end{itemize}
    \item Coefficients $a_{\col'} = \min\{\rhs - RS_r : \text{non-tight } r \text{ with } \btable_{r\col'} = 1\}$
\end{itemize}

\bigskip
\textbf{Example:} Start with $\be{1}{2} + \be{2}{2} + \be{3}{2} \leq 2$
\begin{enumerate}
    \item Lift $\be{1}{4}$ with coefficient 2
    \item Lift $\be{3}{1}$ with coefficient 1
    \item Lift $\be{1}{1}$ with coefficient 1
    \item Result: $\be{1}{1} + \be{1}{2} + 2\be{1}{4} + \be{2}{2} + \be{3}{1} + \be{3}{2} \leq 2$
\end{enumerate}
\end{frame}


\begin{frame}{\lazyCutlift Example: strengthening $\be{1}{2} + \be{2}{2} + \be{3}{2} \leq 2$}


After generating a valid cut, \textbf{strengthen} it by adding more terms

\newcommand{\tight}{\cellcolor{orange!25}}
\let\lifted\tight
\let\blocked\tight
\newcommand{\fracval}{\cellcolor{yellow!25}}
\newcommand{\wholeval}{\cellcolor{blue!15}}
\newcommand{\nonfeasible}{\cellcolor{red!15}}
\newcommand{\inU}{\cellcolor{green!20}}

\centering
\smallskip
\small
\renewcommand{\arraystretch}{1.15}
\begin{tabular}{c|>{\centering}p{\tblcw}>{\centering}p{\tblcw}>{\centering}p{\tblcw}>{\centering}p{\tblcw}|>{\centering}p{\tblcw}>{\centering}p{\tblcw}>{\centering}p{\tblcw}|>{\centering}p{\tblcw}>{\centering}p{\tblcw}>{\centering\arraybackslash}p{\tblcw}|c}
% Column headers: grey background when in Xtight
 & \only<4->{\blocked}$\only<4->{\color{blue}}\besml{1}{1}$
 & \blocked$\color{blue}\besml{1}{2}$
 & \blocked$\besml{1}{3}$
 & \only<2->{\blocked}$\only<2->{\color{blue}}\besml{1}{4}$
 & \blocked$\besml{2}{1}$
 & \blocked$\color{blue}\besml{2}{2}$
 & \only<2->{\blocked}$\besml{2}{3}$
 & \only<3->{\blocked}$\only<3->{\color{blue}}\besml{3}{1}$
 & \blocked$\color{blue}\besml{3}{2}$
 & \only<2->{\blocked}$\besml{3}{3}$
 & $RS$ \\
\hline
% Row 1: RS=1. Tight at step 3 (lift be{3}{1}, a=1, RS->2)
\only<3->{\tight}%
$1$
 & 0 & \tight 1 & 0 & 0
 & 1 & 0 & 0
 & \only<3->{\tight}1 & 0 & 0
 & \only<1>{1}\only<2>{1}\only<3->{\textbf{2}} \\
% Row 2: RS=2. Tight from the start
\tight%
$2$
 & 0 & 0 & 1 & 0
 & 0 & \tight 1 & 0
 & 0 & \tight 1 & 0
 & \textbf{2} \\
% Row 3: RS=0. Tight at step 2 (lift be{1}{4}, a=2, RS->2)
\only<2->{\tight}%
$3$
 & 0 & 0 & 0 & \only<2->{\tight}1
 & 0 & 0 & 1
 & 0 & 0 & 1
 & \only<1>{0}\only<2->{\textbf{2}} \\
% Row 4: RS=1. Tight at step 4 (lift be{1}{1}, a=1, RS->2)
\only<4->{\tight}%
$4$
 & \only<4->{\tight}1 & 0 & 0 & 0
 & 0 & \tight 1 & 0
 & 0 & 0 & 1
 & \only<1-3>{1}\only<4->{\textbf{2}} \\
% Row 5: RS=2. Tight from the start
\tight%
$5$
 & 0 & \tight 1 & 0 & 0
 & 1 & 0 & 0
 & 0 & \tight 1 & 0
 & \textbf{2} \\
\end{tabular}

\bigskip
\only<1>{%
Cut: ${\color{blue}\be{1}{2}} + {\color{blue}\be{2}{2}} + {\color{blue}\be{3}{2}} \leq 2$
\quad {\small Tight rows: \{2,5\}. \colorbox{orange!25}{Tight} columns cannot be lifted.}%
}%
\only<2>{%
Lift ${\color{blue}\be{1}{4}}$ with $a=2$: \quad
${\color{blue}\be{1}{2}} + {\color{blue}2\be{1}{4}} + {\color{blue}\be{2}{2}} + {\color{blue}\be{3}{2}} \leq 2$
\quad {\small row 3 now tight}%
}%
\only<3>{%
Lift ${\color{blue}\be{3}{1}}$ with $a=1$: \quad
${\color{blue}\be{1}{2}} + {\color{blue}2\be{1}{4}} + {\color{blue}\be{2}{2}} + {\color{blue}\be{3}{1}} + {\color{blue}\be{3}{2}} \leq 2$
\quad {\small row 1 now tight}%
}%
\only<4>{%
Lift ${\color{blue}\be{1}{1}}$ with $a=1$: \quad
${\color{blue}\be{1}{1}} + {\color{blue}\be{1}{2}} + {\color{blue}2\be{1}{4}} + {\color{blue}\be{2}{2}} + {\color{blue}\be{3}{1}} + {\color{blue}\be{3}{2}} \leq 2$
\quad \checkmark {\small all tight, all \colorbox{orange!25}{blocked}}%
}%
\end{frame}

}

\ignore{

\begin{frame}{Fractional Cuts (\lazyFractional)}
\begin{columns}[T]
\begin{column}{0.62\textwidth}
\begin{itemize}
    \item \gls{ilp} solvers can call callbacks on \textbf{fractional \gls{lp} solutions} too
    \item Identify fractional columns not appearing in any feasible row
    \item If such a column exists $\Rightarrow$ use it to seed \lazyGenerate
    \item If not ($U = \emptyset$) $\Rightarrow$ solution may be in the convex hull, skip
\end{itemize}

\medskip
\textbf{Example:} $\hat{v} = (0, 0.5, 0.5, 0, 0, 1, 0, 0.5, 0.5, 0)$
\begin{itemize}
    \item Fractional: $F = \{2,3,8,9\}$, whole: $W = \{6\}$
    \item Only feasible row: $D = \{2\}$
    \item Cols 2, 8 not in $D$ $\Rightarrow$ $U = \{2, 8\}$
    \item Cut: $\be{1}{2} + \be{2}{2} \leq 1$
\end{itemize}
\end{column}
\begin{column}{0.35\textwidth}
\centering
\begin{tikzpicture}[scale=0.95]
    % Axes
    \draw[->, gray] (0.5,0.5) -- (3.5,0.5) node[right, font=\scriptsize] {$\x{2}$};
    \draw[->, gray] (0.5,0.5) -- (0.5,3.5) node[above, font=\scriptsize] {$\x{3}$};
    % Ticks
    \foreach \v in {1,2,3} {
        \draw[gray] (\v,0.5) -- (\v,0.42) node[below, font=\tiny] {\v};
        \draw[gray] (0.5,\v) -- (0.42,\v) node[left, font=\tiny] {\v};
    }
    % Convex hull of table rows in (x2, x3) projection
    % Vertices: r1(1,1), r3(3,3), r4(2,3), r5(1,2); r2(2,2) on edge
    \fill[blue!8] (1,1) -- (3,3) -- (2,3) -- (1,2) -- cycle;
    \draw[blue!50, thick] (1,1) -- (3,3) -- (2,3) -- (1,2) -- cycle;
    % Table rows (integer points)
    \fill[blue!80] (1,1) circle (2.5pt);
    \fill[blue!80] (2,2) circle (2.5pt);
    \fill[blue!80] (3,3) circle (2.5pt);
    \fill[blue!80] (2,3) circle (2.5pt);
    \fill[blue!80] (1,2) circle (2.5pt);
    % Row labels
    \node[blue!80, font=\scriptsize, below] at (1,1) {$r_1$};
    \node[blue!80, font=\scriptsize, below right] at (2,2) {$r_2$};
    \node[blue!80, font=\scriptsize, right] at (3,3) {$r_3$};
    \node[blue!80, font=\scriptsize, above] at (2,3) {$r_4$};
    \node[blue!80, font=\scriptsize, left] at (1,2) {$r_5$};
    % Fractional point on hull edge (r1--r5) -- no cut possible
    \fill[orange] (1,1.5) circle (2.5pt);
    \node[orange, font=\scriptsize, left] at (1,1.5) {$\hat{v}'$};
    % Fractional point outside hull -- can be cut
    \fill[red!70!black] (2,1.5) circle (2.5pt);
    \node[red!70!black, font=\scriptsize, right] at (2,1.5) {$\hat{v}$};
    % Separating cut (the hull edge x3=x2 extended)
    \draw[red!70!black, dashed, thick] (0.7,0.7) -- (3.3,3.3) node[font=\scriptsize, above left] {cut};
\end{tikzpicture}

\smallskip
{\scriptsize
\begin{tabular}{c|cc}
 & $\x{2}$ & $\x{3}$ \\ \hline
$r_1$ & 1 & 1 \\
$r_2$ & 2 & 2 \\
$r_3$ & 3 & 3 \\
$r_4$ & 2 & 3 \\
$r_5$ & 1 & 2 \\
\end{tabular}
}

\smallskip
{\scriptsize $\hat{v}'$ is in the hull; $\hat{v}$ is not.}
\end{column}
\end{columns}
\end{frame}

\begin{frame}{\lazyFractional Example: $\hat{v} = (0, 0.5, 0.5, 0, 0, 1, 0, 0.5, 0.5, 0)$}

\centering
\smallskip
\small
\renewcommand{\arraystretch}{1.15}
\begin{tabular}{c|>{\centering}p{\tblcw}>{\centering}p{\tblcw}>{\centering}p{\tblcw}>{\centering}p{\tblcw}|>{\centering}p{\tblcw}>{\centering}p{\tblcw}>{\centering}p{\tblcw}|>{\centering}p{\tblcw}>{\centering}p{\tblcw}>{\centering\arraybackslash}p{\tblcw}}
% Column headers: U columns get blue text from step 3
 & $\besml{1}{1}$
 & $\only<3->{\color{blue}}\besml{1}{2}$
 & $\besml{1}{3}$
 & $\besml{1}{4}$
 & $\besml{2}{1}$
 & $\besml{2}{2}$
 & $\besml{2}{3}$
 & $\only<3->{\color{blue}}\besml{3}{1}$
 & $\besml{3}{2}$
 & $\besml{3}{3}$ \\
\hline
% Row 1: 1-cols = {2,5,8}. col 5 not in F∪W → not in D
\only<2->{\nonfeasible}%
$1$
 & 0 & 1 & 0 & 0
 & \only<2->{\nonfeasible}1 & 0 & 0
 & 1 & 0 & 0 \\
% Row 2: 1-cols = {3,6,9}. All in F∪W → D
\only<2->{\tight}%
$2$
 & 0 & \only<3->{\inU}0 & 1 & 0
 & 0 & 1 & 0
 & \only<3->{\inU}0 & 1 & 0 \\
% Row 3: 1-cols = {4,7,10}. col 4 not in F∪W → not in D
\only<2->{\nonfeasible}%
$3$
 & 0 & 0 & 0 & \only<2->{\nonfeasible}1
 & 0 & 0 & \only<2->{\nonfeasible}1
 & 0 & 0 & \only<2->{\nonfeasible}1 \\
% Row 4: 1-cols = {1,6,10}. col 1 not in F∪W → not in D
\only<2->{\nonfeasible}%
$4$
 & \only<2->{\nonfeasible}1 & 0 & 0 & 0
 & 0 & 1 & 0
 & 0 & 0 & \only<2->{\nonfeasible}1 \\
% Row 5: 1-cols = {2,5,9}. col 5 not in F∪W → not in D
\only<2->{\nonfeasible}%
$5$
 & 0 & 1 & 0 & 0
 & \only<2->{\nonfeasible}1 & 0 & 0
 & 0 & 1 & 0 \\
\hline
% v row: show F (yellow) and W (blue)
$\hat{v}$
 & 0 & \fracval 0.5 & \fracval 0.5 & 0
 & 0 & \wholeval 1 & 0
 & \fracval 0.5 & \fracval 0.5 & 0 \\
\end{tabular}

\bigskip
\only<1>{%
Fractional columns $F = \{$\colorbox{yellow!25}{2,3,8,9}$\}$, whole columns $W = \{$\colorbox{blue!15}{6}$\}$%
}%
\only<2>{%
Feasible rows $D = \{$\colorbox{orange!25}{2}$\}$: all 1-cols $\subseteq F \cup W$.
\quad \colorbox{red!15}{Others} have 1s in zero-columns of $\hat{v}$.%
}%
\only<3>{%
$U = \{$\colorbox{green!20}{2, 8}$\}$: $F$-columns with \colorbox{green!20}{0} in all $D$ rows.
Choose ${\color{blue}\be{1}{2}}$, expand $\Rightarrow$ cut: ${\color{blue}\be{1}{2} + \be{2}{2}} \leq 1$%
}%
\end{frame}

\begin{frame}{Convex Hull of the Running Example}
\centering
\tdplotsetmaincoords{65}{130}
\begin{tikzpicture}[tdplot_main_coords, scale=1.6]
    % Define points
    \coordinate (r1) at (2, 1, 1);
    \coordinate (r2) at (3, 2, 2);
    \coordinate (r3) at (4, 3, 3);
    \coordinate (r4) at (1, 2, 3);
    \coordinate (r5) at (2, 1, 2);
    % Fractional points
    \coordinate (vp) at (2, 1, 1.5);  % v' = midpoint r1,r5 (in hull)
    \coordinate (v)  at (2.5, 2, 1.5); % v  = outside hull

    % Short axes near origin
    \draw[->, gray!50] (0.3,0.3,0.3) -- (2.2,0.3,0.3) node[anchor=north east, font=\scriptsize] {$\x{1}$};
    \draw[->, gray!50] (0.3,0.3,0.3) -- (0.3,1.8,0.3) node[anchor=north west, font=\scriptsize] {$\x{2}$};
    \draw[->, gray!50] (0.3,0.3,0.3) -- (0.3,0.3,2.0) node[anchor=south, font=\scriptsize] {$\x{3}$};
    % Faces of tetrahedron (r1, r3, r4, r5) with transparency
    \fill[blue!12, opacity=0.4] (r1) -- (r4) -- (r5) -- cycle;
    \fill[blue!12, opacity=0.4] (r3) -- (r4) -- (r5) -- cycle;
    \fill[blue!12, opacity=0.5] (r1) -- (r3) -- (r5) -- cycle;
    \fill[blue!12, opacity=0.5] (r1) -- (r3) -- (r4) -- cycle;

    % Edges of tetrahedron
    \draw[blue!40, thick] (r1) -- (r3);
    \draw[blue!40, thick] (r1) -- (r4);
    \draw[blue!40, thick] (r1) -- (r5);
    \draw[blue!40, thick] (r3) -- (r4);
    \draw[blue!40, thick] (r3) -- (r5);
    \draw[blue!40, thick] (r4) -- (r5);

    % Table row points (hull vertices)
    \foreach \p in {(r1),(r3),(r4),(r5)} {
        \fill[blue!80] \p circle (2pt);
    }
    % r2 on edge r1--r3 (not a vertex)
    \fill[blue!50] (r2) circle (2pt);

    % Row labels
    \node[blue!80, font=\scriptsize, below left]  at (r1) {$r_1\,(2,1,1)$};
    \node[blue!50, font=\scriptsize, below right] at (r2) {$r_2\,(3,2,2)$};
    \node[blue!80, font=\scriptsize, right]        at (r3) {$r_3\,(4,3,3)$};
    \node[blue!80, font=\scriptsize, above left]   at (r4) {$r_4\,(1,2,3)$};
    \node[blue!80, font=\scriptsize, left]         at (r5) {$r_5\,(2,1,2)$};

    % Fractional point in hull (on edge r1--r5)
    \fill[orange] (vp) circle (2pt);
    \node[orange, font=\scriptsize, left] at (vp) {$\hat{v}'$};

    % Fractional point outside hull
    \fill[red!70!black] (v) circle (2pt);
    \node[red!70!black, font=\scriptsize, right] at (v) {$\hat{v}$};
\end{tikzpicture}

\vspace{-2mm}
{\scriptsize Tetrahedron: $r_2$ lies on edge $r_1$--$r_3$. \quad
$\hat{v}'$: in hull (no cut). \quad $\hat{v}$: outside hull (can be cut).}
\end{frame}
}





\ignore{

\begin{frame}{Hybrid Eager-Lazy Approach (\lazyCutoff)}
\begin{itemize}
    \item For each table constraint, choose:
    \begin{itemize}
        \item \textbf{Eager} encoding (\gls{mdd}-based) for \emph{small} tables
        \item \textbf{Lazy} cut generation for \emph{large} tables
    \end{itemize}
    \item Threshold based on number of rows (e.g., 1000, 2000, 3000)
    \item Rationale: large tables are costly to encode, and many rows may be irrelevant
\end{itemize}

\bigskip
\begin{center}
\begin{tikzpicture}[>=Stealth, font=\small]
\node[draw, rounded corners, fill=blue!10, minimum width=3cm] (small) at (0,0) {Small tables};
\node[draw, rounded corners, fill=orange!10, minimum width=3cm] (large) at (6,0) {Large tables};
\node[draw, rounded corners, fill=blue!30] (enc) at (0,-1.2) {\baseMddDomIncr};
\node[draw, rounded corners, fill=orange!30] (lazy) at (6,-1.2) {\lazyAll};
\draw[->, thick] (small) -- (enc);
\draw[->, thick] (large) -- (lazy);
\node[draw, rounded corners, fill=green!20, minimum width=4cm] (hybrid) at (3,-2.5) {\lazyCutoffT{threshold}};
\draw[->, thick] (enc) -- (hybrid);
\draw[->, thick] (lazy) -- (hybrid);
\end{tikzpicture}
\end{center}
\end{frame}

	}

\begin{frame}{Results -- CSP (180 instances)}
\centering
\small
\resulttabular{%
\resultencoding{}{}{}
\hline
\resultlazy
{\alt<3->{}{\color{white}}}%
{\alt<4->{}{\color{white}}}%
{\alt<5->{}{\color{white}}}%
}

\smallskip
{\scriptsize pst = posted, sat = feasible, sol = solved. Time = PAR2 avg [s.]}
\end{frame}

\begin{frame}{Eager vs.\ Lazy -- Scatter Plots}

\centering
\begin{figure}
\centering
\includesvg[width=\textwidth,height=0.78\textheight,keepaspectratio]{analysis/scatter-CSP-\compare.svg}
	\caption{\small Solve time [s.]: Eager (top-left) -- Lazy (bottom-right). Light colour = big table}
\end{figure}
\end{frame}

\begin{frame}{Results -- CSP (180 instances)}
\centering
\small
\resulttabular{%
\resultencoding{}{}{}
\hline
\resultlazy{}{}{}
\hline
\resulthybrid
}

\smallskip
{\scriptsize pst = posted, sat = feasible, sol = solved. Time = PAR2 avg [s.]}
\end{frame}

%===============================================================================
\section{Conclusion}
%===============================================================================

\begin{frame}{Conclusion \& Future Work}
\begin{columns}[T]
\begin{column}{0.72\textwidth}
\textbf{Summary:}
\begin{itemize}
    \item Studied eager encodings and lazy cuts for table constraints
    \item Novel MDD-based encodings outperform standard methods
    \item Cutlifting and fractional cuts improve cut generation
    \item Hybrid approach combines the strengths of both paradigms
\end{itemize}

\bigskip
\textbf{Future work:}
\begin{itemize}
    \item Better understanding of which tables suit which method (\eg~sparsity)
    \item Extend to \emph{negative, smart, and reified tables}
\end{itemize}
\end{column}
\begin{column}{0.22\textwidth}
\centering
\includesvg[width=1.6cm]{qr-paper}\\[-2pt]
{\scriptsize Paper}\\[6pt]
\includesvg[width=1.6cm]{qr-code}\\[-2pt]
{\scriptsize Code}\\[6pt]
\includesvg[width=1.6cm]{qr-website}\\[-2pt]
{\scriptsize Website}
\end{column}
\end{columns}
\end{frame}

\begin{frame}[allowframebreaks]{References}
\tiny
\bibliographystyle{plainurl}
\bibliography{references}
\end{frame}

\end{document}
